\documentclass[12pt,a4paper]{article}

% Paquets de langue et d'encodage
\usepackage[utf8]{inputenc}
\usepackage[T1]{fontenc}
\usepackage[french]{babel}

% Paquets pour la mise en page et les mathématiques
\usepackage{geometry}
\geometry{hmargin=2.5cm,vmargin=2.5cm}
\usepackage{amsmath}
\usepackage{amsfonts}
\usepackage{amssymb}
\usepackage{graphicx}
\usepackage{hyperref}
\usepackage{xcolor}
\usepackage{array}
\usepackage{booktabs}
\usepackage{longtable}
\usepackage{multirow}
\usepackage{tabularx}
\usepackage{array}

% Configuration des liens
\hypersetup{
    colorlinks=true,
    linkcolor=blue,
    filecolor=magenta,      
    urlcolor=cyan,
    pdftitle={Rapport : Architecture du Projet MultiAgentMADDPG DevicePowerEnv},
}

\title{\textbf{Rapport Technique} \\ \Large Architecture du Projet : \texttt{MultiAgentMADDPG\_DevicePowerEnv}}
\author{Documentation de Projet (PFE)}
\date{\today}

\begin{document}

\maketitle

\vspace{1cm}

\begin{abstract}
Ce rapport détaille la nouvelle architecture de l'environnement d'Apprentissage par Renforcement Multi-Agents (MARL) baptisé \texttt{MultiAgentMADDPG\_DevicePowerEnv}. Il présente les fondements de la modélisation des objets connectés (IoT), la gestion rigoureuse des contraintes physiques, la dynamique d'exécution concurrente (priorité et files d'attente), le modèle énergétique, ainsi que la chronologie de décision de l'intelligence artificielle lors du choix de l'équipement.
\end{abstract}

\newpage
\tableofcontents
\newpage

\section{Architecture du Projet (Modélisation de l'Environnement)}

Le projet s'appuie sur une architecture d'Apprentissage par Renforcement Multi-Agents (MARL) déployée avec l'algorithme \textbf{MADDPG} (Multi-Agent Deep Deterministic Policy Gradient). Il s'inscrit dans un paradigme d'entraînement centralisé avec exécution décentralisée (CTDE). La dynamique de simulation repose sur un environnement compatible avec l'API \texttt{Gymnasium}, nommé \texttt{MultiAgentIoTEnvLatencyEnergySum}.

\subsection{Entités Principales}
\begin{itemize}
    \item \textbf{Les Agents :} Chaque agent représente une tâche d'intelligence artificielle (par exemple, un réseau de neurones divé en couches). À chaque étape de décision, un agent doit allouer la couche courante à un équipement du réseau (nœud IoT).
    \item \textbf{Le Gestionnaire de Ressources (\texttt{ResourceManager}) :} Il représente le moteur physique de l'environnement. Il valide en temps réel l'état du matériel (CPU, Ram, QoS, Niveau d'exposition et Sécurité) pour chaque nœud avant d'approuver une allocation.
\end{itemize}

\subsection{Espace d'Observation et Fonction de Récompense}
L'objectif du système s'articule autour d'une optimisation multi-objective couplant latence et énergie. La récompense pour une allocation réussie est donnée par :
\begin{equation}
    R_{succ\grave{e}s} = -(T_{\text{total}} + E_{\text{co\hat{u}t}})
\end{equation}
Afin de favoriser un comportement coopératif entre les agents, la récompense finale est une pondération : 70\% de la récompense individuelle et 30\% de la moyenne globale de tous les  agents.

\section{Gestion des Contraintes (Hard Constraints)}

L'environnement intègre un ensemble de règles strictes, appelées \textit{Hard Constraints}, qui garantissent le respect des limitations physiques, sécuritaires et fonctionnelles du système IoT. Toute violation de ces contraintes entraîne un échec immédiat de l'agent, avec une pénalité sévère (récompense de $-500.0$) et la terminaison de l'épisode (\textit{stall termination}).

Ces contraintes traduisent les objectifs fondamentaux du système, qui vise à assurer une allocation intelligente des tâches tout en respectant plusieurs critères critiques :

\begin{enumerate}

    \item \textbf{Contrainte de Latence et Ressources :}  
    Chaque tâche doit être exécutée dans un délai acceptable tout en respectant les capacités matérielles du dispositif cible. Cela inclut :
    \begin{itemize}
        \item Capacité CPU disponible
        \item Mémoire RAM suffisante
        \item Bande passante réseau
    \end{itemize}
    Toute allocation dépassant ces ressources est interdite.

    \item \textbf{Contrainte Énergétique :}  
    L'énergie requise pour exécuter la tâche doit être strictement inférieure à l'énergie disponible dans la batterie du dispositif :
    \begin{equation}
        E_{\text{task}} \leq E_{\text{remaining}}
    \end{equation}
    Cette contrainte intègre un modèle énergétique dynamique basé sur le DVFS (Dynamic Voltage and Frequency Scaling).

    \item \textbf{Contrainte de Confidentialité (Privacy Constraint) :}  
    Chaque couche du modèle possède un niveau de sensibilité (\texttt{privacy\_level}).  
    Le dispositif cible doit posséder un niveau d'autorisation suffisant :
    \begin{equation}
        \texttt{privacy\_clearance} \geq \texttt{privacy\_level}
    \end{equation}
    Sinon, l'exécution est strictement interdite.

    \item \textbf{Contrainte de Confiance (Trust Constraint) :}  
    Chaque dispositif possède un score de confiance dynamique :
    \begin{equation}
        \texttt{trust\_score} \in [0,1]
    \end{equation}
    Si ce score est inférieur à un seuil minimal, l'appareil est considéré comme non fiable et ne peut pas exécuter la tâche.

    \item \textbf{Contrainte de Sécurité et d'Exposition du Modèle :}  
    Afin d'éviter les attaques de type \textit{white-box} et la reconstruction du modèle, le système impose une limite sur la fraction d'exposition du réseau de neurones :
    \begin{equation}
        S_l \leq S_{\text{max}}
    \end{equation}
    où $S_l$ représente la fraction du modèle exécutée par un même dispositif.

    De plus, une contrainte de diversité est imposée :
    \begin{itemize}
        \item Un même nœud ne peut pas exécuter plusieurs couches consécutives
    \end{itemize}

\end{enumerate}

Ces contraintes rendent le problème d'optimisation hautement complexe et non linéaire, justifiant l'utilisation de méthodes d'apprentissage par renforcement multi-agents (MARL) pour apprendre des stratégies optimales dans un environnement dynamique et incertain.

\section{Exécution des Tâches et Modèle de Concurrence}

Dans notre environnement simulé, la latence d'exécution d'une tâche est calculée en prenant en compte à la fois le temps de service et le temps d'attente dans les files de traitement (queues).

\subsection{Modèle de Latency}

\subsubsection{Latence de Calcul}

Le temps de calcul est donné par :

\begin{equation}
    T_{comp}^{service} = \frac{C}{f}
\end{equation}

où :
\begin{itemize}
    \item $C$ est la demande en calcul (computation\_demand)
    \item $f$ est la vitesse du processeur (cpu\_speed)
\end{itemize}

En tenant compte de la file d'attente :

\begin{equation}
    T_{comp} = T_{comp}^{service} + T_{comp}^{wait}
\end{equation}

avec :
\begin{equation}
    T_{comp}^{wait} = \text{device\_comp\_queue}
\end{equation}

 

\subsubsection{Latence de Communication}

La latence de communication dépend du transfert de données entre deux dispositifs.

\begin{equation}
    T_{comm}^{service} = \frac{D}{R}
\end{equation}

où :
\begin{itemize}
    \item $D$ est la taille des données (output\_data\_size)
    \item $R$ est la bande passante (bandwidth)
\end{itemize}

La latence totale de communication devient :

\begin{equation}
    T_{comm} = T_{comm}^{service} + T_{comm}^{wait}
\end{equation}

avec :
\begin{equation}
    T_{comm}^{wait} = \text{device\_comm\_queue}
\end{equation}

 

\subsubsection{Cas Spécifiques de Transmission}

\begin{itemize}
    \item \textbf{Exécution locale :}  
    Si la tâche est exécutée sur le même dispositif que la précédente :
    \begin{equation}
        T_{comm} = 0
    \end{equation}

    \item \textbf{Offloading entre dispositifs :}  
    Si le dispositif courant est différent du précédent :
    \begin{equation}
        D = \text{output\_data\_size}
    \end{equation}

    \item \textbf{Première couche :}  
    Une taille de données initiale constante est utilisée :
    \begin{equation}
        D = D_{init}
    \end{equation}
\end{itemize}

 

\subsubsection{Latence Totale}

La latence totale d'une tâche est donnée par :

\begin{equation}
    T_{total} = T_{comp} + T_{comm}
\end{equation}

\begin{equation}
    T_{total} = (T_{comp}^{service} + T_{comp}^{wait}) + (T_{comm}^{service} + T_{comm}^{wait})
\end{equation}

 

\subsubsection{Mise à Jour des Files d’Attente}

Après allocation d'une tâche :

\begin{itemize}
    \item La file de calcul est mise à jour :
    \begin{equation}
        \text{device\_comp\_queue} \leftarrow \text{device\_comp\_queue} + T_{comp}^{service}
    \end{equation}

    \item La file de communication est mise à jour :
    \begin{equation}
        \text{device\_comm\_queue} \leftarrow \text{device\_comm\_queue} + T_{comm}^{service}
    \end{equation}
\end{itemize}

 

\subsection{Modèle de Consommation Énergétique}

La consommation énergétique est calculée à partir du modèle simplifié :

\begin{equation}
    E = P \times T
\end{equation}

\subsubsection{Énergie de Calcul}

\begin{equation}
    E_{comp} = P_{comp} \times T_{comp}
\end{equation}

\subsubsection{Énergie de Communication}

\begin{equation}
    E_{comm} = P_{comm} \times T_{comm}
\end{equation}

\subsubsection{Énergie Totale}

\begin{equation}
    E_{total} = E_{comp} + E_{comm}
\end{equation}

 

\subsection{Interprétation}

Ce modèle reflète fidèlement la dynamique réelle du système :

\begin{itemize}
    \item Les files d’attente introduisent de la congestion
    \item L’offloading génère un coût de communication
    \item L’énergie dépend directement du temps d’exécution
\end{itemize}

Ainsi, les agents MARL apprennent à :
\begin{itemize}
    \item éviter les dispositifs surchargés
    \item réduire les communications inutiles
    \item équilibrer charge et performance
\end{itemize}

\section{Profil de Puissance Dynamique (Hardware-Proportional Power)}

À l'inverse des anciens prototypes utilisant des pénalités globales "codées en dur", la puissance de consommation $P$ (en Watts) est calculée de manière proportionnelle pour chaque matériel via mise à l'échelle :
\begin{equation}
    P_{comp, i} = \text{Base}_{\text{Power\_Comp}} \times \left( \frac{\text{CPU\_Speed}_i}{\text{Base}_{\text{CPU\_Speed}}} \right)
\end{equation}
\begin{equation}
    P_{comm, i} = \text{Base}_{\text{Power\_Comm}} \times \left( \frac{\text{Bandwidth}_i}{\text{Base}_{\text{Bandwidth}}} \right)
\end{equation}
L'équation énergétique permettant aux agents MADDPG de calculer l'énergie ponctionnée sur des serveurs hétérogènes devient ainsi une véritable signature mathématique :
\begin{equation}
    E_{\text{co\hat{u}t\_total}} = (P_{\text{comp}, i} \times T_{\text{comp}}) + (P_{\text{comm}, i} \times T_{\text{comm}})
\end{equation}
\textit{Théorème métier :} Un processeur ultra-rapide consomme plus de Watts, mais sa latence diminuée (le multiplicateur Temps) garantira une pondération globale évitant l'explosion des coûts en batterie.

\section{Processus de Décision de l'Agent (Inférence IA)}

Le processus de sélection du nœud IoT cible par l’agent connecte la modélisation métier aux mathématiques du Deep Learning. Ce flux d'exécution séquentiel garantit que l'IA optimise son utilité tout en respectant inconditionnellement les limites physiques de l'environnement matériel.

\subsection{1. Le Filtrage de Légalité (Valid Actions)}
Avant d'invoquer ses réseaux neuronaux, l'agent reçoit de la part de l'environnement une sous-liste appelée \texttt{valid\_actions}. L'environnement pré-trie tous les appareils en excluant d'office ceux ne répondant pas aux \textit{Hard Constraints} (Batterie morte, RAM saturée, ou Score de Confiance insuffisant).

\subsection{2. Dilemme d'Exploration-Exploitation ($\epsilon$-greedy)}
L'agent utilise un générateur stochastique gouverné par le paramètre Epsilon ($\epsilon$). 
\begin{itemize}
    \item \textbf{Exploration} : Avec une probabilité $\epsilon$, l'agent met son intelligence en veille et choisit un appareil purement au hasard parmi les actions valides. Ceci garantit la découverte de nouvelles approches.
    \item \textbf{Exploitation} : Avec une probabilité $1 - \epsilon$, l'agent requiert l'inférence de son modèle cognitif (le réseau \texttt{DiscreteActor}).
\end{itemize}

\subsection{3. Inférence Neuronale et Vecteur d'Intentions (Logits)}
En phase d'exploitation, l'agent transmet le tenseur d'état du système à son réseau Actor. Le réseau neuronal produit en retour un vecteur appelé \textit{Logits}. Ce vecteur contient une note pour chaque périphérique de l'infrastructure. Un logit élevé représente l'intuition de l'IA que cet équipement matériel minimisera conjointement à cet instant T l'équation Latence + Énergie.

\subsection{4. Contrainte par l'Infini (Action Masking)}
Le machine learning est faillible et peut prédire une note très élevée pour un appareil physiquement incapable d'assumer la tâche. Le code contraint mathématiquement la prédiction via un filtre d'élagage :
l'Action Masking.
Tous les appareils ne figurant pas dans \texttt{valid\_actions} se voient assigner de force la note de $-\infty$ (moins l'infini).
Les appareils autorisés par l'environnement conservent leur score généré par le réseau de neurones.

\subsection{5. La Transposition par Argmax (Choix Final)}
En dernier aboutissement, la prise de décision est le résultat de l'extraction \texttt{argmax} sur le vecteur masqué complet. 
Cet opérateur renvoie systématiquement l'appareil légal possédant le score théorique le plus attractif, verrouillant ainsi la pertinence de l'offloading IA par des barrières de sécurité purement programmatiques interdisant tout crash structurel de l'écosystème.

\vspace{1cm}
\begin{center}
    \rule{0.5\textwidth}{0.4pt} \\
    \vspace{0.2cm}
    \textit{Fin du rapport sur l'architecture environnementale.}
\end{center}


\begin{table*}[ht]
\centering
\small
\renewcommand{\arraystretch}{1.3}

\begin{tabularx}{\textwidth}{|p{2.2cm}|X|X|X|X|X|X|}
\hline
\textbf{Article} & \textbf{Problématique} & \textbf{Méthodologie} & \textbf{Modèle/Algorithme} & \textbf{Objectif} & \textbf{Contributions} & \textbf{Limites} \\
\hline

Baccour et al. (2024) &
Privacy + Latence dans inference distribuée IoT &
Formulation optimisation + MARL &
MARL coopératif/compétitif &
Minimiser latence + préserver privacy &
Distribution des CNN sans serveur central, agents autonomes &
Complexité élevée, simulation seulement \\
\hline

RL-PDNN (2021) &
Protection contre attaques (white-box) en DNN distribué &
Optimisation $\rightarrow$ RL (DQN) &
Deep Q-Learning &
Minimiser latence + protéger modèle &
Distribution dynamique des couches CNN + privacy-aware &
Moins adapté aux systèmes multi-agents complexes \\
\hline

Mi et al. (2024) &
Offloading dynamique avec énergie variable &
Modélisation queue + DRL &
MAPPO (MARL) &
Minimiser délai des tâches &
CTDE, D2D + Energy harvesting &
Hypothèses fortes (modèle simplifié) \\
\hline

AI on the Edge (2024) &
Déploiement AI sous contraintes énergie/mémoire &
AutoML + optimisation multi-objectif &
Evolutionary + RL-based design &
Optimiser performance + énergie &
Pareto optimization + prédiction énergétique &
Pas focalisé sur offloading dynamique \\
\hline

MEC Survey &
Limites cloud (latence, bande passante) &
Survey + taxonomy &
N/A &
Comprendre MEC &
Architecture MEC + applications &
Pas de solution algorithmique \\
\hline

Notre Travail (2024/2026) &
Allocation DNN sous contraintes multiples (Énergie, Trust, Sécurité) en environnement hétérogène IoT &
Simulation avancée + MARL (Resource Manager strict/souple) &
MADDPG / MATD3 (CTDE, Actor-Critic) &
Optimiser (Latence + Énergie) avec Privacy et Sécurité ($S_l$) &
Contraintes unifiées (Trust, Privacy Level, exposition $S_l$), intégration DVFS (CPU/GPU) &
Grand espace d'état, convergence lente, tuning nécessaire \\
\hline

\end{tabularx}

\caption{Comparative Literature Matrix for MEC and MARL-based Task Offloading}
\label{tab:literature_matrix}

\end{table*}

\end{document}